%% kapa: exam=1 round=35 year=2024 period=1 form=A subject=economics subject_ko=경제학원론 seq=36 number=76
%% source: Q-Net 감정평가사 기출문제 (cst00309) artlSeq=5213229 file=raw/1008/5213229/1교시 문제지 원본.pdf
%% conversion: kordoc 4.24.51 -> tools/gen1.py (auto); review=proofread-sonnet; flags=none
\begin{question}{76}
아래의 폐쇄경제 IS-LM 모형에서 중앙은행은 균형이자율을 현재보다 5 만큼 높이는 긴축적 통화정책을 실시하여 균형국민소득을 감소시키고자 한다. 현재 명목화폐공급량($M$)이 40일 때, 이를 달성하기 위한 명목화폐공급량의 감소분은? (단, $r$은 이자율, $Y$는 국민소득, $M^{d}$는 명목화폐수요량, $P$는 물가수준이고 1로 고정되어 있다.)
\begin{framedtext}
○ IS 곡선: $r=120-5Y$ \par
○ 실질화폐수요함수: $\frac{M^{d}}{P}=3Y-r$ \par
\end{framedtext}
\begin{choices}
  \item 5
  \item 8
  \item 10
  \item 15
  \item 20
\end{choices}
\end{question}
\answer{2}
